\subsection*{解三角}
\begin{enumerate}
    \item 在 $\bigtriangleup  {ABC}$ 中,角 $A,B,C$ 对应的边分别为 $a,b,c,A = \dfrac{\pi }{6},a = 2$，当 $b$ 为何值时, ${\Delta ABC}$ 的面积 $S$ 最大？并求出 $S$ 的最大值.
    \begin{jieda}
        由余弦定理得 ${a^2} = {b}^{2} + {c}^{2} - {2bc}\cos A$ ,即 $4 + \sqrt{3}{bc} = {b}^{2} + {c}^{2}$ .
        又 ${b}^{2} + {c}^{2} \geq  {2bc}$ ,则 $4 + \sqrt{3}{bc} \geq  {2bc}$ ,
        所以 ${bc} \leq  \dfrac{4}{2 - \sqrt{3}} = 4\left( {2 + \sqrt{3}}\right)$ ,当且仅当 $b = c = \sqrt{2} + \sqrt{6}$ 时,等号成立.
        则 $S = \dfrac{1}{2}{bc}\sin A = \dfrac{1}{4}{bc} \leq  2 + \sqrt{3}$ .
        因此,当 $b$ 为 $\sqrt{2} + \sqrt{6}$ 时, ${S}_{\bigtriangleup {ABC}}$ 取得最大值 $2 + \sqrt{3}$ .
    \end{jieda}


    \item 在 $\bigtriangleup {ABC}$ 中,角 $A,B,C$ 的对边分别是 $a,b,c$ ,若 $a = 2,b = 3,\angle {ACB}$ 的平分线 ${CD}$ 的长为 $\dfrac{4\sqrt{6}}{5}$ ,则 ${AB}$ 边上的中线 ${CE}$ 的长为\blkda{$\dfrac{\sqrt{17}}{2}$}.
    \begin{jieda}
        \begin{dingautolist}{172}
            \item 方法1
            
            根据角分线定理可设$AD =3x$，$BD = 2x$，故根据余弦定理有$\dfrac{4+\dfrac{96}{25}-4x^2}{4 \times \dfrac{4\sqrt{6}}{5}} = \dfrac{9+\dfrac{96}{25}-9x^2}{6 \times \dfrac{4\sqrt{6}}{5}}$，故$12 + \dfrac{288}{25} - 12x^2 = 18 + \dfrac{192}{25} - 18x^2$，故$x = \dfrac{3}{5}$
            \item 方法2

            大三角形面积等于两个小三角形面积，$\dfrac{1}{2} \times \dfrac{4\sqrt{6}}{5} \times 2 \times sin\left(\dfrac{C}{2}\right) + \dfrac{1}{2} \times \dfrac{4\sqrt{6}}{5} \times 3 \times sin\left(\dfrac{C}{2}\right) = \dfrac{1}{2} \times 3 \times 2 \times sinC$，整理得$2sin\dfrac{C}{2}(3cos\dfrac{C}{2}-\sqrt{6}) = 0$，因为$sin\dfrac{C}{2} > 0$，所以$cos\dfrac{C}{2} = \dfrac{\sqrt{6}}{3}$，故$cosC = \dfrac{1}{3}$
        \end{dingautolist}
        故$AB = 3$，$AE = BE = \dfrac{3}{2}$

        设$CE = y$，根据余弦定理有$\left\{\begin{array}{l}
                4 = \dfrac{9}{4} + y^2 - 3ycos\theta  \\
                9 = \dfrac{9}{4} + y^2 + 3ycos\theta
        \end{array}\right.$，两式相加的可得$x = \dfrac{\sqrt{17}}{2}$
    \end{jieda}

    \item  在锐角 $\bigtriangleup  {ABC}$ 中,内角 $A,B,C$ 的对边分别为 $a,b,c$ ,已知 $c \cdot  \cos A = \left( {{2b} - a}\right) \cos C$ .
    \begin{enumerate}
        \item 求 $\angle C$ 的值
        \item 求 $\dfrac{a}{b}$ 的取值范围.
    \end{enumerate}
    \begin{jieda}
        \begin{enumerate}
            \item 由 $\left( {{2b} - a}\right) \cos C = c\cos A$ 及正弦定理得： $\left( {2\sin B - \sin A}\right) \cos C = \sin C\cos A$ .
            
            $\therefore 2\sin B\cos C - \sin A\cos C = \sin C\cos A$ ,可得: $2\sin B\cos C = \sin \left( {A + C}\right)  = \sin B$
            
            $\because C \in  \left( {0,\pi }\right) ,\sin B \neq  0$ ,且 $\bigtriangleup {ABC}$ 是锐角三角形，$\therefore \cos C = \dfrac{1}{2}$ ,可得: $C = \dfrac{\pi }{3}$ .

            \item $\because C = \dfrac{\pi }{3},\therefore B = \dfrac{2\pi }{3} - A$ .
            $\because 0 < A < \dfrac{\pi }{2},\because 0 < \dfrac{2\pi }{3} - B < \dfrac{\pi }{2},\therefore \dfrac{\pi }{6} < B < \dfrac{\pi }{2}$ .
            $\therefore \tan B > \dfrac{\sqrt{3}}{3}$ .
            
            $\dfrac{a}{b} = \dfrac{\sin A}{\sin B} = \dfrac{\sin \left( {\dfrac{2\pi }{3} - B}\right) }{\sin B} = \dfrac{\dfrac{\sqrt{3}}{2}\cos B + \dfrac{1}{2}\sin B}{\sin B} = \dfrac{1}{2} + \dfrac{\sqrt{3}}{2\tan B} \in  \left( {\dfrac{1}{2},2}\right) .$
        \end{enumerate}
    \end{jieda}


    \item 在锐角 $\bigtriangleup  {ABC}$ 中,内角 $A,B,C$ 所对应的边分别是 $a,b,c$ ,且${2c}\sin \left( {B - A}\right)  = {2a}\sin A\cos B + b\sin {2A}$ ,则 $\dfrac{c}{a}$ 的取值范围是\blkda{$(1,2)$}.
    \begin{jieda}
        由正弦定理和正弦二倍角公式可得$2\sin C\sin \left( {B - A}\right)  = 2\sin A\sin A\cos B + \sin B\sin {2A} = 2\sin A\sin A\cos B + 2\sin B\sin A\cos A = 2\sin A\left( {\sin A\cos B + \sin B\cos A}\right) = 2\sin A\sin \left( {A + B}\right)$，故 $\sin \left( {B - A}\right)  = \sin A$
        
        因为 $0 < A < \dfrac{\pi }{2},0 < B < \dfrac{\pi }{2}$ ,所以 $- \dfrac{\pi }{2} < B - A < \dfrac{\pi }{2}$，所以 $B = {2A},C = \pi  - {3A}$ 

        由正弦定理得 $\dfrac{c}{a} = \dfrac{\sin C}{\sin A} = \dfrac{\sin {3A}}{\sin A} = 3-4sin^2A$ .
        
        由 $0 < B = {2A} < \dfrac{\pi }{2},0 < C = \pi  - {3A} < \dfrac{\pi }{2}$ 可得 $\dfrac{\pi }{6} < A < \dfrac{\pi }{4}$，故$\dfrac{c}{a} \in \left( {1,2}\right)$
    \end{jieda}
    
\end{enumerate}